% Source: https://artofproblemsolving.com/community/c3h1883598_quick_observation
% Stable identity: aops-topic-1883598-q01
The radius $r$ satisfies $\pi r^2=16\pi$, so $r=4$. Join the centre $O$ to the endpoints of one side: the triangle has sides $4$, $4$, $4$, so it is equilateral and the side subtends a central angle of $60^\circ$.

\begin{asy}
usepackage("amsmath");
size(4.5cm);
defaultpen(fontsize(10pt));
draw(circle((0,0),4));
pair[] V;for(int k=0;k<6;++k) V.push(4*dir(60*k));
draw(V[0]--V[1]--V[2]--V[3]--V[4]--V[5]--cycle,linewidth(1));
fill((0,0)--V[0]--V[1]--cycle,rgb(0.85,0.92,1));
draw((0,0)--V[0]--V[1]--cycle);
dot((0,0));label("$O$",(0,0),SW);
label("$4$",V[0]/2,S);label("$4$",V[1]/2,W);label("$4$",(V[0]+V[1])/2,SW);
draw(arc((0,0),0.9,0,60));label("$60^\circ$",1.45*dir(30),fontsize(8pt));
label("\mbox{}",(min(currentpicture,user=true)+max(currentpicture,user=true))/2,black); // keeps the text after the figure black when LaTeX typesets the labels inline
\end{asy}

All sides of a regular polygon subtend equal central angles that together make $360^\circ$, so the polygon has $\dfrac{360}{60}=\boxed{6}$ sides.
